\section{A biquadratic form with no dominated square}
\label{sec:form}

The geometric argument will produce bilinear forms whose squares are
bounded by the Choi--Lam form $b$ in~\eqref{eq:form}.
We prove directly the particular consequence of its classical
extremality~\cite[Theorem~4.4]{choi-lam-1977} that the geometric
argument needs: every such bilinear form must vanish. This property is
also called \emph{weak extremality}~\cite[Definition~2.7]{blekherman-raita-shankar-sinn-2022}.

\begin{lemma}\label{lem:form}
The form $b:\R^3\times\R^3\to\R$ defined in~\eqref{eq:form}
is nonnegative and nonzero. If a real bilinear form
$\ell:\R^3\times\R^3\to\R$ satisfies
\[
 \ell(z,y)^2\leq b(z,y)
 \qquad\text{for every }z,y\in\R^3,
\]
then $\ell=0$.
\end{lemma}

\begin{proof}
Put $a_i=z_i^2$, with indices read cyclically modulo $3$.
The matrix of $b(z,y)$ as a quadratic form in $y$ is
\[
 M(z)=
 \begin{pmatrix}
 a_1+a_3&-z_1z_2&-z_1z_3\\
 -z_1z_2&a_2+a_1&-z_2z_3\\
 -z_1z_3&-z_2z_3&a_3+a_2
 \end{pmatrix}.
\]
Its diagonal entries are nonnegative. Its three principal minors of
order two are
\[
 \det M(z)[\{i,i+1\}]
 =a_i^2+a_{i-1}(a_i+a_{i+1})\geq0,
 \qquad i=1,2,3,
\]
where the brackets denote a principal submatrix. Finally,
\[
 \det M(z)
 =a_1^2a_2+a_2^2a_3+a_3^2a_1-3a_1a_2a_3\geq0
\]
by the arithmetic--geometric mean inequality. All principal minors
are therefore nonnegative, so $M(z)\succeq0$ and $b(z,y)\geq0$.
For the standard coordinate vectors $e_1,e_2,e_3$ of $\R^3$,
we have $b(e_1,e_1)=1$.

Now write $\ell(z,y)=\sum_{i,j}c_{ij}z_i y_j$ and assume the stated
bound. Zeros will eliminate the off-diagonal coefficients; fourth-order
vanishing along curves will eliminate the diagonal coefficients.
At every zero of $b$, the form $\ell$ vanishes.
In particular,
\[
 b(e_i,e_{i-1})=0
 \quad\Longrightarrow\quad c_{i,i-1}=0.
\]
For each sign vector $\sigma\in\{-1,1\}^3$, we also have
$b(\sigma,\sigma)=3-6+3=0$, and hence
\[
 0=\sum_i c_{ii}
   +\sum_{i<j}(c_{ij}+c_{ji})\sigma_i\sigma_j.
\]
Multiplying by $\sigma_p\sigma_q$ and averaging over the eight sign
vectors gives $c_{pq}+c_{qp}=0$ for every $p<q$.
Together with $c_{i,i-1}=0$, these identities force all off-diagonal
entries to vanish. Thus $\ell(z,y)=\sum_i c_i z_i y_i$ for some real $c_i$.

For a cyclic index $i$ and $s\in\R$, set
\[
 z=e_i+s e_{i-1},\qquad y=e_{i-1}+s e_i.
\]
Direct substitution gives
\[
 b(z,y)=2s^2-2s^2+s^4=s^4,
 \qquad \ell(z,y)=s(c_i+c_{i-1}).
\]
For $s\ne0$, the bound implies $(c_i+c_{i-1})^2\leq s^2$.
Letting $s\to0$ yields $c_i+c_{i-1}=0$ for all three indices.
These equations force $c_1=c_2=c_3=0$, proving the lemma.
\end{proof}
